% Chapter 2: Writing applications with QueenMQ
% The chapter's own opening text goes after \chapter; its sections are the files below,
% in chapters/2-writing-applications/, in the order they appear in the book.

\chapter{Writing applications with QueenMQ}

For this book, we chose to show code with the QueenMQ JavaScript SDK, because we think it is one of the easiest languages for understanding how the system works. On the website, you will find all the examples of this book in every QueenMQ SDK.

% Chapter 2, section 1: Setup a local broker
% Write in plain LaTeX. Code of a guide step: \guide{<step>}; commands: \begin{shell}...\end{shell};
% a screenshot: \screenshot{file}{caption}; a gap in the draft: \placeholder{...}.
% See book/latex/queen.sty for the full list.

\section{Setup a local broker}

\screenshot[b]{dashboard-overview}{The dashboard, served on the same port as the API.}

One binary serves the HTTP API and the dashboard on one port:

\begin{shell}
docker run -p 6632:6632 ghcr.io/queen-mq/queen:2.7.0
\end{shell}

Open http://localhost:6632 and watch every step of this part land in the dashboard
(Figure~\ref{fig:dashboard-overview}). Then add the client:

\queenfig{one-node}{One process serves the API, the dashboard, health and metrics on port 6632.}

\begin{shell}
npm install queen-mq
\end{shell}

% Chapter 2, section 2: Create and configure queues
% Write in plain LaTeX. Code of a guide step: \guide{<step>}; commands: \begin{shell}...\end{shell};
% a screenshot: \screenshot{file}{caption}; a gap in the draft: \placeholder{...}.
% See book/latex/queen.sty for the full list.

\section{Create and configure queues}

A queue is a named set of partitions that share one configuration. You never have to create one:
the first push that names a queue creates it, with the defaults. Configure it when you want other
settings.

\queenfig{queue-anatomy}{A queue: options shared by its partitions, each an ordered list of messages.}

\guide{configure}

\why{a queue created by a push is good enough to start, and most queues never need more. When they
do, the settings live on the queue, in one place, instead of in every worker.}

\subsection{Every option}

Here is a queue that sets most of them, and the table below says what each one does.

\guide{configure-jobs}

\begin{options}{Option}{Default}{What it does}
leaseTime & 60 s (300 s from \texttt{create()}) & How long a worker holds the messages it took before they go to another worker. \\
retryLimit & 3 & How many failed attempts a message gets. After the last one it goes to the dead-letter queue. \\
retryDelay & 1000 ms & How long a failed message waits before it comes back. Its partition waits with it, for that consumer group. \\
deadLetterQueue, dlqAfterMaxRetries & true & File a message that ran out of retries in the dead-letter queue. With both off, it is dropped. \\
delayedProcessing & 0 s & How long every new message waits after its push before anyone can take it. \\
windowBuffer & 0 s & How long new messages wait in their partition, so a burst is handed over as one batch. \\
dedupWindowSeconds & 3600 s & How long a partition remembers \texttt{transactionId}s. \texttt{0} turns deduplication off. \\
retentionEnabled & false & Must be true for the two options below to do anything. \\
retentionSeconds & 0 & Remove messages older than this, consumed or not. \\
completedRetentionSeconds & 0 & Remove messages older than this that every consumer group has already acked. \\
maxWaitTimeSeconds & 0 & Remove messages older than this, consumed or not. Works without \texttt{retentionEnabled}. \\
encryptionEnabled & false & Encrypt payloads on disk. Every node needs \texttt{QUEEN\_ENCRYPTION\_KEY}. \\
namespace, task & from the name & Labels that group queues in the dashboard and in a discovery pop. \\
\end{options}

\paragraph{Leases.} A worker that takes messages holds them under a lease. Until it acks them or the
lease runs out, no other worker of the same consumer group gets that partition.

\why{a worker can crash in the middle of its work. The lease is how its messages get back to
someone else without a human noticing. Set it longer than your slowest handler, or renew it while
the handler runs (see Pop messages).}

\paragraph{Retries.} A failed message comes back \texttt{retryLimit} times, each time after
\texttt{retryDelay}, and then goes to the dead-letter queue.

\why{most failures are temporary: a database that restarts, a provider that times out. Trying again
later fixes them. Waiting between tries gives the other side time to recover, and a message that
can never succeed stops after a few tries instead of blocking its partition forever.}

\paragraph{Delays.} \texttt{delayedProcessing} holds every new message of the queue for a while.
\texttt{windowBuffer} holds a group's new messages until the burst is over, so the handler gets them
together.

\why{some work should not start at once (a reminder sent five minutes after a signup), and some
work is cheaper in bulk (one database write for a burst of fifty updates). For a delay that is
different for each message, use a timer instead (see KV, Locks and Timers).}

\paragraph{Deduplication.} Inside the dedup window, a second push with the same
\texttt{transactionId} in the same partition is not written.

\why{producers retry when they do not get an answer, and a retry must not create a second order.}

\paragraph{Retention.} By default Queen keeps every message, and the disk grows. Retention runs only
when \texttt{retentionEnabled} is true; then \texttt{retentionSeconds} removes old messages, read or
not, and \texttt{completedRetentionSeconds} removes old messages that every group has finished.
\texttt{maxWaitTimeSeconds} removes old messages without needing \texttt{retentionEnabled}. A node
whose disk is 85\% full refuses the writes that would grow it.

\why{an acked message is not deleted: it stays in its partition so that a new consumer group, an
audit or a fixed bug can read it again. Retention is how you decide how long that history is worth
its disk.}

\paragraph{Encryption.} With \texttt{encryptionEnabled}, payloads are encrypted on disk.

\why{some payloads hold personal data that must not sit on a disk in clear text.}

\subsection{Three traps}

\begin{itemize}
\item \texttt{config(...).create()} in the JavaScript client sends your options together with the
client's own defaults (a 300-second lease, encryption off, no retention). Calling it on an existing
queue resets every option you leave out. Name every option you care about.
\item On a node without \texttt{QUEEN\_ENCRYPTION\_KEY}, an encrypted queue stores plain text, and
nothing warns you.
\item \texttt{priority}, \texttt{ttl}, \texttt{maxSize} and \texttt{minPopWaitTime} are accepted and
stored, and do nothing.
\end{itemize}

On the site, with the code in all six SDK languages: \href{https://queenmq.com/writing-applications/queues/}{queenmq.com/writing-applications/queues}.

% Chapter 2, section 3: Push messages
% Write in plain LaTeX. Code of a guide step: \guide{<step>}; commands: \begin{shell}...\end{shell};
% a screenshot: \screenshot{file}{caption}; a gap in the draft: \placeholder{...}.
% See book/latex/queen.sty for the full list.

\section{Push messages}

In Kafka you choose a number of partitions and hash keys onto them, so a slow customer holds up
every customer that shares its partition. Here a partition is one entity: a customer, an order, a
conversation. The push that names it creates it, order holds inside it, and one queue has held ten
million.

The \texttt{transactionId} makes a retry safe: the same id within the dedup window (an hour by default) is
not written twice.

\guide{push}

\why{order matters per entity (an order's events must be handled in sequence), not across the
whole queue. One partition per entity keeps that order without one slow entity holding up the
others. A push with no partition goes to a partition called \texttt{Default}.}

\subsection{Every message gets its own answer}

One push can carry many messages, and the broker answers each one: \texttt{queued} (stored, with
its offset), \texttt{duplicate} (that \texttt{transactionId} is already there, nothing was written)
or \texttt{error} (refused, nothing written). The JavaScript client throws when any message is
refused, unless you handle it yourself, and never throws for a duplicate.

\queenfig{push-verdicts}{One push, one answer per message.}

\guide{push-callbacks}

\why{a duplicate is usually good news: it means an earlier attempt already got through. A refusal
is not. Telling them apart per message lets one bad item in a batch of a hundred be retried alone,
and lets a producer that retried after a timeout know that its first attempt worked.}

\subsection{Buffering, for many small messages}

Each push is one HTTP request. A producer that pushes thousands of small messages a second can
let the client collect them and send them in batches instead.

\guide{push-buffer}

A buffer is sent when it holds \texttt{messageCount} messages or when \texttt{timeMillis} has
passed, whichever comes first. If the producer gets far ahead of the network (four times
\texttt{messageCount} waiting), a buffered push waits instead of letting memory grow.

\queenfig{push-buffer}{Client buffering: many small pushes, a few requests.}

\why{one request per message spends most of its time on the round trip, not on the message.
Batching gets far more messages through the same connection. The price: a message is stored only
when its batch is sent, so a process that crashes loses what was still in its buffer. Call
\texttt{flushAllBuffers()} before you exit, and do not buffer messages you cannot afford to lose.}

\subsection{When several messages must land together}

A push of several messages is not all-or-nothing: each message gets its own answer, so some can be
stored while others are refused. When messages must be stored together or not at all (an order and
its payment request, say), push them in a transaction (see Advanced transactional usage).

\subsection{Binary push, compression and framed pop}

By default the client and the broker speak JSON. Three options of the client change how messages
travel, and nothing else: the calls, their results, deduplication and the guarantees stay the
same.

\begin{options}{Option}{Default}{What it does}
binaryPush & false & A push, direct or buffered, goes as one binary body instead of a JSON document. A transaction that pushes sends its pushes the same way. \\
pushCompression & false & Needs \texttt{binaryPush}. The client packs each partition's batch the way the broker stores it, and compresses it with zstd once. The broker keeps those bytes as they are. \\
framedPop & false & \texttt{pop()} and \texttt{consume()} get the records as they are stored, and the client unpacks them into the same messages. \\
\end{options}

Turn them on when you create the client:

\guide{wire-forms}

zstd comes with Node.js, so there is nothing to install. Some requests stay JSON by themselves: a
push with a \texttt{traceId}, a pop of a conflating group, and a pop of an encrypted queue. Acks,
KV, timers and locks do not change.

\why{at a high message rate, the broker spends CPU turning JSON into stored records
and back, and the network carries every message uncompressed. These options move that work into
the clients, where many processes share it, and send the batches compressed. When neither the
broker's CPU nor the network is the limit, leave them off: JSON is easier to read when you debug.}

One trap: \texttt{binaryPush} and \texttt{pushCompression} need broker 2.3.0 or later on every node.
An older broker answers 400, and the client does not fall back to JSON, so turn them on after the
last node is upgraded. \texttt{framedPop} is safe with any broker: an older one ignores it and
answers JSON.

On the site, with the code in all six SDK languages: \href{https://queenmq.com/writing-applications/push/}{queenmq.com/writing-applications/push}. More there: \href{https://queenmq.com/writing-applications/clients/}{the binary forms in every SDK}.

% Chapter 2, section 4: Pop messages
% Write in plain LaTeX. Code of a guide step: \guide{<step>}; commands: \begin{shell}...\end{shell};
% a screenshot: \screenshot{file}{caption}; a gap in the draft: \placeholder{...}.
% See book/latex/queen.sty for the full list.

\section{Pop messages}

There are two ways to read. \texttt{consume()} is a loop: it takes messages, calls your handler,
acks them when the handler returns, and goes on. \texttt{pop()} is one call: it hands you the
messages and leaves the ack to you.

\why{most workers are loops, and \texttt{consume()} writes the loop for you, with the error handling
done right. \texttt{pop()} is for the cases a loop does not fit: a scheduled job that reads once, a
script, or code that wants to decide itself when to ack.}

\queenfig{consume-loop}{What consume() does for you.}

\subsection{Queue mode and consumer groups}

Without a consumer group, each message goes to one worker, and the broker sizes every pop for you.

\guide{consume}

With a consumer group, every group gets every message and the workers of one group share them. A
new group starts at the tail, so this one asks for everything already in the queue. Here every
option is set by hand, and the handler acks.

\guide{consume-group}

Ack \texttt{failed} and the message comes back, up to the queue's \texttt{retryLimit}; after the last retry it
goes to the dead-letter queue.

\why{one queue often feeds several services: billing, emails and analytics all want every order.
Each service is a consumer group, reading at its own pace, and none of them can slow the others
down. Inside a group, the workers compete, so you add workers to go faster.}

\queenfig{queue-mode-vs-groups}{Queue mode shares the messages; every consumer group gets all of them.}

\subsection{One message at a time, or a whole batch}

With \texttt{.each()}, the handler gets one message at a time. Without it, the handler gets the
whole batch as an array.

\guide{consume-batch}

What the client does with your handler's result (\texttt{autoAck}, on by default):

\begin{itemize}
\item \textbf{It returns:} the message (or the whole batch) is acked.
\item \textbf{It throws:} the message (or the whole batch) is nacked as \texttt{failed}, which
spends one retry, and the loop goes on with the next messages.
\item \textbf{It throws under \texttt{.each()}:} the rest of that partition's messages in the same
batch are skipped, because they will be delivered again after the failed one, in order. Messages of
other partitions in the batch are handled as usual.
\item \textbf{With \texttt{.autoAck(false)}:} the handler acks itself (in a transaction, for
example). A handler that throws is still nacked, so its partition is not stuck until the lease
runs out.
\end{itemize}

\why{a batch is cheaper (one ack for a hundred messages, one database write for a hundred rows),
and one at a time is simpler and fails one message instead of a hundred. Pick per worker.}

If you add \texttt{.onSuccess()} or \texttt{.onError()} to a consumer, the client stops acking and
nacking for you: you decide in those callbacks what happens to each message.

\guide{consume-callbacks}

\why{some failures should be retried and some never will succeed (an email address that does not
exist). The callbacks let you send the first kind back with \texttt{failed} and the second kind
straight to the dead-letter queue with \texttt{dlq}.}

\subsection{How many messages a pop takes: autopilot}

By default the broker decides how many messages each pop gets (\texttt{batch}, from 100 to 1,000,
based on how fast the group acks) and from how many partitions (\texttt{partitions}, the ones that
have work, shared among the waiting workers). A value you set yourself is never changed.
\texttt{partitions(N)} lets one pop take messages from up to \texttt{N} partitions (1 to 64), with
\texttt{batch} as the total. \texttt{.autopilot(false)} goes back to one message from one partition
per pop.

\why{the right batch size depends on your handler and your load, and it changes during the day.
The broker sees both, so it can size every pop better than a number fixed in the code.}

\subsection{Leases and long work}

A pop leases its messages to the worker for the queue's \texttt{leaseTime}. Work that can take
longer renews the lease while it runs: \texttt{.renewLease(true, 10000)} renews every 10 seconds,
and \texttt{queen.renew(message)} renews once. In PHP, the renewal happens between messages, so a
single long handler calls the renew callback it is given. A lease that has already
expired cannot be renewed: its messages go to the next worker, marked as one delivery attempt
more.

\why{the lease is what protects order and stops two workers from doing the same work. A short lease
gets a crashed worker's messages moving again quickly; renewing lets a slow but healthy worker keep
them.}

\subsection{Where a new group starts}

A group's starting point is decided on its first pop and then stored: \texttt{new} (the default)
reads only what is pushed after that first pop, \texttt{all} reads everything still kept, and
\texttt{subscriptionFrom} reads from a point in time. A group that already exists moves with a seek,
to a time or to the end.

\queenfig{subscription-start}{Where a new group starts, fixed at its first pop.}

\guide{seek}

\why{a new service usually wants only new events, a new projection wants the whole history, and a
bug fix wants to replay the last day. A seek to the end is how a group skips a backlog it no longer
needs.}

\subsection{Only the newest message: conflation}

A group with \texttt{.conflation()} gets only the newest message of each partition, and its cursor
moves past the older ones.

\guide{conflation}

\why{some queues carry commands where only the latest one matters: ``recompute the price of A-1''.
A worker that is behind should do the work once, with the newest input, instead of three times with
inputs that are already old.} Conflation is decided by the group's first pop, needs a consumer
group, and cannot be combined with \texttt{commitOnDelivery}.

\subsection{Faster, with no ack: commitOnDelivery}

A pop with \texttt{.commitOnDelivery()} moves the group's cursor as the messages are handed out.
There is no lease and nothing to ack.

\guide{commit-on-delivery}

\why{a live dashboard or a metrics chart can lose a reading now and then, and would rather not pay
for leases and acks. This is at-most-once delivery: a worker that crashes after the pop loses that
batch. It works with \texttt{pop()} only; \texttt{consume()} refuses it.}

\subsection{Running a consumer for real}

\guide{consume-run}

\begin{itemize}
\item \texttt{concurrency(n)} runs \texttt{n} handlers at once, each on its own partitions.
\item \texttt{limit(n)} stops after \texttt{n} messages, and \texttt{idleMillis(ms)} stops after
that long with nothing to do.
\item \texttt{wait(true)} (the default) waits for messages instead of returning empty, up to
\texttt{timeoutMillis(ms)} (30 seconds by default, at most 60).
\item A \texttt{signal} stops the loop at once. Messages it had taken but not started go back with a
\texttt{retry} ack, which spends no retry, so a deploy does not leave partitions waiting for their
leases to run out.
\end{itemize}

\why{a worker runs for weeks, through deploys and quiet nights. These options decide how much it
does at once, when it stops, and how it gives its work back when it is told to.}

\subsection{Read it again}

Consuming does not delete: the queue is a log, kept for its retention. A new group can start from a
point in time and read the same orders again, while the groups that already read them stay where
they are.

\guide{replay}

Or read a partition by offset with \texttt{fetch}: no group, no lease, no ack. It suits an audit,
rebuilding one entity, or a sink that keeps its own position.

\guide{fetch}

\why{history is useful long after the first reader is done with it: to build a new view, to audit
what happened, or to rebuild one customer's state from its events.}

On the site, with the code in all six SDK languages: \href{https://queenmq.com/writing-applications/pop/}{queenmq.com/writing-applications/pop}. More there: \href{https://queenmq.com/writing-applications/clients/}{reads by offset in every SDK}.

% Chapter 2, section 5: Acknowledge messages
% Write in plain LaTeX. Code of a guide step: \guide{<step>}; commands: \begin{shell}...\end{shell};
% a screenshot: \screenshot{file}{caption}; a gap in the draft: \placeholder{...}.
% See book/latex/queen.sty for the full list.

\section{Acknowledge messages}

An ack tells the broker what happened to a message. A \texttt{completed} ack moves the consumer
group's cursor past it, and that is all: the message stays in its partition until retention removes
it.

\queenfig{cursor-per-group}{An ack moves the group's cursor; the message stays.}

\why{the cursor is the group's memory of where it is. Keeping the message after the ack is what
lets other groups, and later replays, read it too.}

\subsection{Four answers}

\begin{options}{Status}{Spends a retry}{What happens, and when to use it}
completed & no & Done. The cursor moves past it. \texttt{true} means the same. \\
failed & yes & The lease is released and the message comes back after \texttt{retryDelay}. After the last retry it goes to the dead-letter queue. Use it for failures that may pass. \texttt{false} means the same. \\
retry & no & The lease is released and the message comes back at once. Use it to give work back (a worker that is shutting down). \\
dlq & no & Filed in the dead-letter queue now. Use it when trying again cannot help (a payload that will never parse). \\
\end{options}

A \texttt{failed} or \texttt{dlq} ack can carry the reason (\texttt{error:~'card~declined'} next to
the group), and the dead-letter queue keeps it.

\why{``it did not work'' has different meanings. Sometimes trying again in a while is right,
sometimes now, and sometimes never. Each answer sends the message where it should go.}

\subsection{One message or a whole batch}

\texttt{queen.ack()} takes one message or an array. With an array, one call acks the whole batch.

\guide{pop-ack}

Each message of an array can carry its own answer:

\guide{ack-each-status}

In a batch, the cursor stops at the first message that will come back: a \texttt{completed}
message after it is delivered again too.

\queenfig{batch-ack-rule}{In a batch, the cursor stops at the first message that comes back.}

\why{one call for a hundred messages is much cheaper than a hundred calls. The cursor rule keeps
order: a message is never passed while an earlier one of the same partition is still to be done.}

\subsection{Three things to know}

\begin{itemize}
\item \textbf{Inside \texttt{.each()}, throw to fail a message.} When the handler throws, the client
nacks it and skips the rest of that partition's batch. When the handler acks \texttt{failed} itself
and returns, the client does not know: it goes on with the partition's next messages, their acks are
refused (that part of the batch is coming back anyway), and they are handled twice.
\item \textbf{Every ack answers HTTP 200, even a refused one.} Read \texttt{success} in the answer: a
refused ack usually means the lease ran out and another worker has the message now.
\item \textbf{An ack carries its consumer group and its lease.} The client takes both from the
message. Over HTTP you send them yourself: an ack without its lease is not checked, and can complete
a batch another worker holds.
\end{itemize}

On the site, with the code in all six SDK languages: \href{https://queenmq.com/writing-applications/ack/}{queenmq.com/writing-applications/ack}.

% Chapter 2, section 6: Retries and the dead-letter queue
% Write in plain LaTeX. Code of a guide step: \guide{<step>}; commands: \begin{shell}...\end{shell};
% a screenshot: \screenshot{file}{caption}; a gap in the draft: \placeholder{...}.
% See book/latex/queen.sty for the full list.

\section{Retries and the dead-letter queue}

Every \texttt{failed} ack spends one retry. The message comes back after the queue's
\texttt{retryDelay}, and its partition waits with it, for that consumer group. After
\texttt{retryLimit} failures, the message goes to the dead-letter queue and the group moves on. You
set both on the queue (see Create and configure queues).

\why{retrying is how temporary failures heal by themselves, and the limit is how a message that can
never succeed stops blocking everything behind it in its partition.}

Two cases do not spend a retry: a \texttt{retry} ack, and a lease that runs out. So a message that
crashes its worker's process every time comes back forever and never reaches the dead-letter queue.
If some input can crash your worker, check for it before the work starts and ack such a message with
\texttt{dlq}.

\queenfig{retry-dlq}{The life of a message: four endings, and the replay.}

\subsection{What a dead letter keeps}

A dead letter keeps the payload, the error the worker sent, the number of retries, and the consumer
group that gave up on it. It stays until you replay it or delete it, or until its queue holds too
many: past 100,000 dead letters, 1 GiB or 14 days in one queue, the oldest go first (from broker
2.7.0; each cap is a setting, and \texttt{0} turns it off). Each group has its own: a message that
failed for billing can be fine for analytics.

\why{a failure you cannot see is a failure you cannot fix. The dead-letter queue is where they wait
for a person, or for a fix in the code.}

\subsection{Read it}

\guide{dlq-read}

\texttt{limit()} and \texttt{offset()} page through it, newest first. The client also has
\texttt{from()} and \texttt{to()}, but the 2.7.0 broker ignores them.

\subsection{Replay it}

Once the cause is fixed, put the messages back:

\guide{dlq-replay}

A replay moves the dead letter back into its partition, at the end, as a new message with the
\texttt{transactionId} \texttt{dlq:<id>}, and removes it from the dead-letter queue in the same
step. Doing it twice cannot create a second copy. Every group that reads that partition sees the
replayed message, not only the group that gave up on it. Over HTTP the same thing is
\texttt{POST /api/v1/dlq/:id/replay}.

\why{the fix usually comes after the failure: a bug deployed, a provider back up. Replay lets the
same messages go through again without anyone pushing them by hand.}

\subsection{Skip it, or turn it off}

An ack with \texttt{dlq} files a message straight away, without spending its retries. Setting
\texttt{deadLetterQueue} and \texttt{dlqAfterMaxRetries} to false on the queue drops a message
that runs out of retries instead of keeping it.

\why{some messages are worth nothing once they fail (a typing indicator, a stale price), and
keeping them would only fill a list nobody reads.}

On the site, with the code in all six SDK languages: \href{https://queenmq.com/writing-applications/retries/}{queenmq.com/writing-applications/retries}. More there: \href{https://queenmq.com/writing-applications/pop/}{consuming}.

% Chapter 2, section 7: Advanced transactional usage
% Write in plain LaTeX. Code of a guide step: \guide{<step>}; commands: \begin{shell}...\end{shell};
% a screenshot: \screenshot{file}{caption}; a gap in the draft: \placeholder{...}.
% See book/latex/queen.sty for the full list.

\section{Advanced transactional usage}

The ack, the order's new state, the next event and tomorrow's reminder commit as one log entry, or
none of them does. That one call replaces the outbox table, the idempotency table and the cron job.

\guide{transaction}

Then open the dashboard: its KV page holds the order's new state, and its Timers page holds
tomorrow's reminder.

\why{a worker's step is several writes: it marks its input done, saves what changed, tells the
next service and plans what comes later. When those writes go to different systems, a crash between
two of them leaves the step half done, and someone writes code to find and repair the halves. One
commit has no halves.}

If the worker took too long and its lease ran out, the whole transaction is refused, so a slow
worker can never write over the one that has the message now. If the step committed and only the
answer got lost, a retry with the same \texttt{transactionId}s is refused whole, and nothing is
written twice.

\queenfig{lease-fence}{A worker whose lease ran out: its whole transaction rolls back.}

On the site, with the code in all six SDK languages: \href{https://queenmq.com/writing-applications/transactions/}{queenmq.com/writing-applications/transactions}.

% Chapter 2, section 8: KV, Locks and Timers
% Write in plain LaTeX. Code of a guide step: \guide{<step>}; commands: \begin{shell}...\end{shell};
% a screenshot: \screenshot{file}{caption}; a gap in the draft: \placeholder{...}.
% See book/latex/queen.sty for the full list.

\section{KV, Locks and Timers}

The transaction above already used two of the three: \texttt{kv.put} stored the order's state, and
\texttt{timer()} scheduled tomorrow's reminder, which arrives as a message.

\why{KV is there so workers keep nothing between messages: the state of each order lives next to
its events and changes in the same commit. Timers are there because ``do this in 24 hours'' is a
message that should arrive later, in the order's own partition, not a cron job that scans a table.
Scheduling the same timer key again moves it, so a reminder can be pushed back or cancelled.}

KV also works on its own, outside a transaction. Every write names its lifetime, and the version a
read returns makes the next write conditional:

\guide{kv}

\why{two workers can read the same cart and both write it. With \texttt{expect}, the second write
finds the version moved and is refused, with the current value in the answer, so nobody overwrites
what they never saw. A counter with \texttt{max} is a rate limit in one call.}

The third is a lock. Some work no message triggers: a nightly report, a cleanup, a singleton. Take a
lock for it. A lock is a lease, and a slow holder can outlive it, so guard the transaction with it:
the push commits only while the lock is still yours.

\guide{lock}

\why{a cron job running on three replicas runs three times. A lock makes it run once, and the guard
makes sure a holder that froze and lost its lock cannot commit after another one took over. A
semaphore (\texttt{queen.semaphore(name, n)}) allows \texttt{n} holders at once.}

\queenfig{lock-fence}{A lock expires without telling its holder; the guard keeps it out.}

On the site, with the code in all six SDK languages: \href{https://queenmq.com/writing-applications/kv-locks-timers/}{queenmq.com/writing-applications/kv-locks-timers}.

% Chapter 2, section 9: Ephemeral queues
% Write in plain LaTeX. Code of a guide step: \guide{<step>}; commands: \begin{shell}...\end{shell};
% a screenshot: \screenshot{file}{caption}; a gap in the draft: \placeholder{...}.
% See book/latex/queen.sty for the full list.

\section{Ephemeral queues}

Ephemeral queues live in memory, outside the log, for request/reply. The worker answers on whatever
inbox a request names; the caller makes a private inbox and waits on it. There is no transaction or
replay here, and a crash loses what is in flight.

The worker:

\guide{worker}

The caller:

\guide{caller}

\why{a request that waits for an answer (``what is the price of A-1, now?'') needs speed more
than history. Writing it to disk, replicating it and keeping it for days would cost time for
nothing. Ephemeral queues keep the same groups, leases and acks, and skip the log. Use a regular
queue for anything that must survive a crash.}

\queenfig{ephemeral-reply}{Request and reply over ephemeral queues, with a private inbox.}

On the site, with the code in all six SDK languages: \href{https://queenmq.com/writing-applications/ephemeral/}{queenmq.com/writing-applications/ephemeral}.

% Chapter 2, section 10: Stream processing
% Write in plain LaTeX. Code of a guide step: \guide{<step>}; commands: \begin{shell}...\end{shell};
% a screenshot: \screenshot{file}{caption}; a gap in the draft: \placeholder{...}.
% See book/latex/queen.sty for the full list.

\section{Stream processing}

A stream reads a queue and writes its results to another. The partition is the key, so each
customer gets its own windows. Like a new group it starts at the tail: push a few more orders and a
result arrives a minute later.

\guide{stream}

A stream reads its queue as a consumer group (\texttt{streams.<queryId>}) and keeps its state, the
open windows, in the broker. Each time it closes windows, it writes the new state, pushes the
results and acks the messages they came from in one transaction.

\queenfig{stream-cycle}{One cycle of a stream is one transaction.}

\why{counting, summing and grouping events over time is a job most systems hand to a separate
stream processor, with its own cluster and its own idea of what was processed. Here it runs in a
normal process, its state lives next to the data, and one commit keeps the results and the inputs in
step: a crash can neither lose a result nor count an event twice.}

\subsection{Four kinds of window}

\begin{options}{Window}{Example}{What it is, and when to use it}
windowTumbling & \texttt{seconds: 60} & Fixed windows, one after the other. Revenue per minute. \\
windowSliding & \texttt{size: 600, slide: 60} & Overlapping windows: the last 10 minutes, every minute. Each event counts in several windows, so it keeps more state. \\
windowSession & \texttt{gap: 1800} & One window per burst of activity, closed after a quiet gap. A visit, a conversation. \\
windowCron & \texttt{every: 'hour'} & Windows aligned to the clock (\texttt{second}, \texttt{minute}, \texttt{hour}, \texttt{day}, \texttt{week}). Hourly and daily reports. \\
\end{options}

\queenfig{window-shapes}{The four kinds of window on one partition's events.}

\guide{stream-sliding}

\guide{stream-session}

\why{``in the last ten minutes'', ``per visit'' and ``per hour'' are different questions, and each
window shape answers one of them without code of your own.}

\subsection{When it happened, or when it arrived}

By default a window holds the events that arrived during it. Give \texttt{eventTime} a function that
reads the time from the message, and windows hold the events that happened during it instead. It
gets the whole message, so the event's own fields are under \texttt{data}; the aggregate's functions
get the event itself.

\guide{stream-hourly}

\texttt{allowedLateness} says how far behind the newest event an event may be and still count. A
later one is dropped, or with \texttt{onLate: 'include'} it opens a fresh window, so the window is
sent a second time with only the late events. \texttt{gracePeriod} keeps a window open a little
after its end, and \texttt{idleFlushMs} (5 seconds) closes the windows of partitions that went quiet.
\texttt{subscriptionMode: 'all'} makes a new stream read what the queue already holds.

\why{events arrive late and out of order: a phone that was offline, a batch upload, a retry. An
hourly report should count a sale in the hour it happened, not in the hour its message got through.}

\subsection{Shaping the data}

\begin{itemize}
\item \texttt{map(fn)} changes each message, \texttt{filter(fn)} drops the ones you do not want,
and \texttt{flatMap(fn)} turns one message into several values.
\item \texttt{aggregate(\{ ... \})} folds a window into named numbers. The name decides the fold:
\texttt{count}, \texttt{sum}, \texttt{min}, \texttt{max} and \texttt{avg}; any other name adds up.
\texttt{reduce(fn, initial)} is the general form, for anything else.
\item \texttt{keyBy(fn)} groups by something other than the partition (per page instead of per user).
Several workers can then write the same key, which slows them down; when that matters, push to a
queue partitioned by the new key first, and stream from it.
\end{itemize}

\why{most stream jobs are the same few steps: drop noise, reshape, group, fold. Having them as
building blocks keeps a stream to a few readable lines.}

\subsection{Two ways to end a stream}

\texttt{.to(queue)} pushes each result to another queue, in the same commit as the state and the
acks, into the partition of the same name unless you pass \texttt{\{ partition \}}.
\texttt{.foreach(fn)} calls your function for each result, and acks only after it returns: a crash
in the middle calls it again.

\why{a result that stays inside Queen can be exactly once. A call to the outside world (an email,
an API) cannot, so \texttt{foreach} makes it at least once and says so. For an external effect that
must happen once, send the result to a queue and let a worker do it with an idempotency key.}

\subsection{Rate limits with a gate}

\texttt{gate(fn)} decides, message by message, whether a message may go on now, with state kept per
partition. Two ready-made gates are included: \texttt{tokenBucketGate} and
\texttt{slidingWindowGate}.

\guide{stream-gate}

A message the gate refuses stops its partition there: it and the ones after it are delivered again,
in order, when the lease runs out.

\why{an SMS provider, a partner API or a rate-limited email service accepts only so many calls a
second. A gate per customer keeps you inside those limits without losing order and without a queue
of your own for the waiting messages.}

\subsection{Running a stream}

\texttt{run()} takes the stream's \texttt{queryId}, its identity: its state is stored under it, so a
process restarted with the same id carries on with the same windows, and several processes with the
same id share the partitions between them. \texttt{batchSize} (200) and \texttt{maxPartitions} (4)
size each cycle, and \texttt{stop()} on the handle \texttt{run()} returns stops it cleanly.

\why{a stream is a long-running worker like any other: it has to survive restarts, scale out, and
stop during a deploy without losing its place.}

On the site, with the code in all six SDK languages: \href{https://queenmq.com/writing-applications/streams/}{queenmq.com/writing-applications/streams}.

Push, consume, ack, replay, fetch, transactions, KV, timers, locks, streams and ephemeral queues all
combine. Build on them and Queen carries the logic you would otherwise write around a broker.

